% Generated by paper/Makefile from metadata.yaml (abstract).
\newcommand{\PreprintAbstract}{Kim and Thorne (1991) computed the renormalized stress-energy of a conformally coupled massless scalar field near the Cauchy horizon of a two-mouth wormhole time machine and found that it diverges as the inverse cube of the distance to the horizon. This paper replicates the computational part of that work: the closed-form results for the covering-space geometry, the defocusing factors, the stress-energy at the throat and between the mouths, the polarized hypersurfaces and the four-momentum are recomputed by an independent numerical route from the article's own starting expressions, with arbitrary-precision arithmetic, and compared with the printed result. All of them are reproduced except one geometric relation (eq. 31), and the stress-energy at the throat (eq. 51) only with the transverse unit vector read as defined in eq. 39. At the throat the residual is identified analytically with a term the closed form does not contain; between the mouths, off the axis, it scales like the correction the text says it ignores; for the other leading-order results it falls with the small parameter, while the exact geometric relations agree to the working precision. Four inconsistencies in the printed text are documented. First, the point-splitting operator as printed (eq. 4a) lacks a factor 1/2 relative to the convention that reproduces the Stefan-Boltzmann density of a thermal state, the primary arbiter; the printed results carry that normalization consistently, so they are reproduced by the operator as printed and lie a factor of two above the arbiter's convention. Second and third, eqs. 31 and 52 are inconsistent with their own context. Fourth, one order-of-magnitude value after eq. 14 does not follow from the inputs printed beside it. None of these changes the article's qualitative conclusions. Every computed number in this paper carries a machine-checked origin (cited, derived, computed or heuristic), and a test forbids any random number generator in the package.}
